% !TeX root=../main.tex

\chapter{مروری بر ادبیات موضوع}
\label{ch:litreview}

\section{مقدمه}
این پایان‌نامه در تقاطع چهار حوزهٔ پژوهشی قرار دارد که تا پیش از این، عمدتاً به‌طور مستقل از یکدیگر پیش رفته‌اند: (۱) بازنمایی معنا و سنجش شباهت معنایی در مدل‌های زبانی، (۲) یادگیری تقابلی به‌عنوان چارچوبی برای شکل‌دادن به فضاهای بازنمایی، (۳) تفسیرپذیری مکانیستی و به‌طور خاص خودرمزگذارهای پراکنده به‌عنوان ابزار گشودن برهم‌نهی ویژگی‌ها، و (۴) ریزتنظیم کارآمد پارامتر و روش‌های حفظ رفتار مدل در حین آن. این فصل هر یک از این چهار حوزه را به‌ترتیب و با تمرکز بر آن دسته از پژوهش‌هایی مرور می‌کند که مستقیماً پایه‌های مفهومی، روش‌شناختی یا فنی فصل سوم را می‌سازند، و در پایان (بخش~\ref{sec:lit-summary}) نشان می‌دهد این پایان‌نامه دقیقاً در کدام نقطهٔ تلاقی این خطوط پژوهشی جای می‌گیرد و چه شکافی را هدف قرار داده است.

\section{بازنمایی و شباهت معنایی در مدل‌های زبانی}
\label{sec:lit-representation}

\subsection{از بازنمایی‌های ایستا تا بازنمایی‌های بافتاری}
پیش از ظهور معماری ترنسفورمر، رایج‌ترین رویکرد برای بازنمایی معنای واژه‌ها، آموختن یک بردار \emph{ایستا}\LTRfootnote{Static}
برای هر واژه بود که مستقل از بافت جمله، ثابت می‌ماند؛ روش
\lr{word2vec}~\cite{mikolov2013word2vec}
نمونهٔ برجستهٔ این رویکرد است که با پیش‌بینی واژه‌های همسایه در یک پنجرهٔ متنی (یا برعکس، پیش‌بینی یک واژه از روی همسایگانش)، بردارهایی می‌آموزد که روابط معنایی ساده--مانند قیاس‌های معنایی که با عملیات برداری ساده (جمع و تفریق) قابل‌بازتولیدند--را تا حدی در هندسهٔ فضای برداری بازتاب می‌دهند. محدودیت اصلی این رویکرد، ناتوانی در نمایش چندمعنایی\LTRfootnote{Polysemy}
و وابستگی به بافت بود: یک واژهٔ یکسان، در جمله‌های مختلف، معنای متفاوتی می‌تواند داشته باشد (برای نمونه، یک واژه که هم به‌معنای عضوی از بدن و هم به‌معنای یک ابزار به کار می‌رود)، اما بازنمایی ایستا--چون تنها یک بردار برای هر واژه در نظر می‌گیرد، صرف‌نظر از جمله‌ای که آن واژه در آن ظاهر شده--این تفاوت را نادیده می‌گیرد و ناچار میانگینی از تمام معانی ممکن آن واژه را در یک بردار فشرده می‌کند.

معرفی معماری ترنسفورمر~\cite{vaswani2017attention}
و مدل‌های پیش‌آموزش‌دیدهٔ مبتنی بر آن--به‌ویژه
\lr{BERT}~\cite{devlin2019bert}
در سمت رمزگذارهای دوسویه و معماری‌های رمزگشا-محور یک‌سویه در سمت دیگر--این محدودیت را با تولید بازنمایی‌هایی \emph{بافتاری}\LTRfootnote{Contextual}
برطرف کرد: بازنمایی هر واژه، تابعی از کل جمله‌ای (یا در مورد مدل‌های یک‌سویه، کل متن پیش از آن) است که در آن ظاهر شده است. مکانیسم توجه\LTRfootnote{Attention}--هستهٔ محاسباتی ترنسفورمر--این بازنمایی‌های بافتاری را با اجازه‌دادن به هر موقعیت برای «توجه‌کردن» با وزن‌های متفاوت به تمام موقعیت‌های دیگر دنباله می‌سازد؛ این وزن‌ها خود از حاصل‌ضرب داخلی بردارهای «پرسمان» و «کلید» هر موقعیت به دست می‌آیند و بنابراین، برخلاف بازنمایی ایستا، در هر پیش‌روی مجدداً و بر اساس محتوای خاص همان جمله محاسبه می‌شوند. در معماری رمزگشا--که پایهٔ مدل‌های این پایان‌نامه نیز هست--خروجی هر لایه (چه بلوک توجه و چه بلوک پیش‌خور) به‌صورت جمعی، به یک مسیر برداری مشترک موسوم به جریان باقیمانده\LTRfootnote{Residual Stream}
افزوده می‌شود؛ به بیان دیگر، بازنمایی نهایی یک نشانه در عمق لایهٔ
$\ell$، جمع تعبیهٔ اولیهٔ آن نشانه با خروجی تک‌تک لایه‌های پیش از
$\ell$
است. همین جریان باقیمانده است که در سرتاسر این پایان‌نامه به‌عنوان محل استخراج بازنمایی جمله در نظر گرفته شده است، و ماهیت جمعی و انباشتی آن--هر لایه، چیزی به بازنمایی قبلی \emph{می‌افزاید}، نه آنکه آن را از نو بسازد--در بخش~\ref{sec:lit-superposition} اهمیت ویژه‌ای پیدا می‌کند.

با این‌همه، بافتاری‌بودن بازنمایی، به‌خودی‌خود تضمینی برای \emph{معناداری} آن نیست. مدل‌های زبانی بزرگ عمدتاً با هدف پیش‌بینی نشانهٔ بعدی آموزش می‌بینند، هدفی که لزوماً بازنمایی سطح جمله را در جهت سنجه‌های معنایی موردنظر انسان (مانند هم‌ارزی معنایی یا استلزام منطقی) شکل نمی‌دهد؛ بازنمایی حاصل، هر آنچه برای پیش‌بینی آماری واژهٔ بعدی مفید باشد را رمزگذاری می‌کند، که ممکن است شامل نشانه‌های سطحی و آماری زبان (مانند همایندی واژگانی) به همان اندازه‌ای باشد که شامل معنای عمیق‌تر است، بی‌آنکه مدل ابزاری برای تفکیک این دو داشته باشد. این شکاف، انگیزهٔ اصلی خط پژوهشی بعدی است.

\subsection{شفافیت پژوهشی خانوادهٔ مدل پایه}
از آنجا که این پایان‌نامه از یک مدل رمزگشا-محور\LTRfootnote{Decoder-only}
خاص
(\lr{Pythia}) استفاده می‌کند، مروری کوتاه بر انگیزهٔ این انتخاب، مفید است. معماری رمزگشا-فقط--که در آن مدل تنها بر پایهٔ متن پیشین، نشانهٔ بعدی را پیش‌بینی می‌کند--امروزه غالب‌ترین انتخاب برای مدل‌های زبانی عام‌منظوره است، چرا که با تولید متن سازگارتر از معماری دوسویهٔ
\lr{BERT}
است. خانوادهٔ
\lr{Pythia}~\cite{biderman2023pythia}--مدل پایهٔ این پایان‌نامه--از میان گزینه‌های موجود در این خانواده، به‌طور خاص با هدف \emph{شفافیت پژوهشی} طراحی شده است: تمام نقاط بازرسی میانی پیش‌آموزش (نه فقط چک‌پوینت نهایی)، ترتیب دقیق دادهٔ دیده‌شده در هر مرحله، و کد کامل پیش‌آموزش، به‌صورت عمومی منتشر شده‌اند. همین فلسفهٔ باز، زمینه‌ساز پیدایش اکوسیستمی از ابزارهای تحلیلی--از جمله همان خودرمزگذارهای پراکندهٔ از پیش‌آموزش‌دیده‌ای که این پایان‌نامه مستقیماً از آن‌ها استفاده می‌کند (بخش~\ref{sec:sae})--حول این خانواده از مدل‌ها شده است؛ بدون وجود چنین خودرمزگذارهایی به‌صورت آماده و عمومی، روش پیشنهادی این پایان‌نامه به آموزش از پیش یک خودرمزگذار اختصاصی نیاز می‌داشت، مرحله‌ای که خود می‌توانست منبع نامعینی اضافی--مستقل از پرسش اصلی پژوهش--باشد. کوچکی اعضای این خانواده نیز، همان‌طور که در بخش~\ref{sec:base-model} اشاره خواهد شد، اجرای شبکهٔ آزمایشی گستردهٔ این پژوهش را در محدودهٔ منابع محاسباتی در دسترس ممکن ساخت.

\section{تشخیص بازنویسی و سنجش شباهت معنایی متن}
\label{sec:lit-paraphrase}
تشخیص بازنویسی\LTRfootnote{Paraphrase Identification}
و سنجش شباهت معنایی متن\LTRfootnote{Semantic Textual Similarity, STS}
دو وظیفهٔ کلاسیک در پردازش زبان طبیعی‌اند که مستقیماً همین شکاف را می‌آزمایند: در وظیفهٔ نخست، از مدل خواسته می‌شود دو جمله را به‌عنوان بازنویسی یکدیگر یا نه، طبقه‌بندی کند؛ در وظیفهٔ دوم، از مدل خواسته می‌شود یک امتیاز پیوستهٔ شباهت (معمولاً در بازهٔ صفر تا پنج) به هر جفت‌جمله نسبت دهد. رویکرد رایج برای استخراج بازنمایی قابل‌مقایسه از یک جملهٔ کامل--به‌جای یک واژهٔ منفرد--استفاده از معماری دوقلو\LTRfootnote{Siamese}
است، که در آن دو شبکه با وزن‌های مشترک، دو جمله را به‌طور مستقل به دو بردار می‌برند و سپس شباهت این دو بردار (معمولاً با شباهت کسینوسی) سنجیده می‌شود.
\lr{Sentence-BERT}~\cite{reimers2019sbert}
نمونهٔ شناخته‌شدهٔ این رویکرد است که با ریزتنظیم
\lr{BERT}
در یک ساختار دوقلو و بر روی جفت‌جمله‌های برچسب‌خورده، بازنمایی‌هایی می‌آموزد که شباهت کسینوسی‌شان--بدون نیاز به عبور هر دو جمله از مدل به‌طور هم‌زمان، برخلاف رویکرد متقاطع اصلی
\lr{BERT}--با شباهت معنایی مورد نظر انسان هم‌بستگی بالایی داشته باشد؛ این جدایی محاسباتی (رمزگذاری مستقل هر جمله) همچنین پیش‌نیاز فنی رویکردی است که این پایان‌نامه نیز از آن استفاده می‌کند، چرا که در آن نیز هر جمله از یک جفت، به‌طور مستقل از مدل عبور داده می‌شود.

\lr{SimCSE}~\cite{gao2021simcse}
گامی فراتر برداشت و نشان داد حتی بدون برچسب انسانی، صرف عبوردادن یک جمله دو بار از مدل با دو دراپ‌اوت تصادفی متفاوت--و در نظرگرفتن این دو خروجی اندکی متفاوت به‌عنوان یک جفت مثبت در یک تابع هزینهٔ تقابلی (بخش~\ref{sec:lit-contrastive})--به‌تنهایی بازنمایی‌های جمله‌ای به‌مراتب بهتری نسبت به بازنمایی‌های خام
\lr{BERT}
تولید می‌کند؛ در نسخهٔ نظارت‌شدهٔ همین کار، جفت‌های استلزام منطقی برچسب‌خورده به‌عنوان جفت مثبت و جفت‌های تناقض به‌عنوان جفت منفی سخت به کار رفته‌اند. این یافته، اهمیت \emph{صورت‌بندی تابع هزینه} (نه صرفاً معماری یا حجم داده) را برای کیفیت بازنمایی جمله برجسته می‌کند--نکته‌ای که مستقیماً به رویکرد این پایان‌نامه (طراحی یک تابع هزینهٔ تقابلی/ناوردا، هرچند نه برای آموزش از پایه بلکه برای شکل‌دادن دوبارهٔ بازنمایی یک مدل از پیش‌آموزش‌دیده) مرتبط است.

\subsection{PAWS و شکاف میان شباهت لفظی و معنایی}
روش‌هایی مانند
\lr{SBERT}
و
\lr{SimCSE}
عمدتاً بر روی مجموعه‌داده‌هایی نظیر
\lr{STS}
ارزیابی می‌شوند که در آن‌ها شباهت لفظی و شباهت معنایی معمولاً هم‌جهت‌اند: جملاتی که واژگان مشترک کمی دارند، معمولاً از نظر معنایی نیز کم‌شباهت‌اند، و برعکس. این هم‌بستگی، یک راه میان‌بُر بالقوه برای مدل فراهم می‌کند: مدلی که صرفاً همپوشانی واژگانی یا حتی ترتیب کیسه‌ای\LTRfootnote{Bag-of-words}
واژه‌ها را می‌سنجد، می‌تواند بدون فهم واقعی ساختار نحوی یا معنایی، عملکرد نسبتاً خوبی روی چنین مجموعه‌داده‌هایی کسب کند.

مجموعه‌دادهٔ
\lr{PAWS}~\cite{zhang2019paws}
دقیقاً برای شکستن همین هم‌بستگی طراحی شد. روش ساخت آن دو مرحله دارد: نخست، با جابه‌جایی ترتیب کلمات یک جملهٔ اصلی (به‌گونه‌ای که ساختار دستوری حاصل هنوز قابل‌قبول باشد)، جمله‌های جدیدی تولید می‌شود که همپوشانی واژگانی بسیار بالایی با جملهٔ اصلی دارند اما بسته به اینکه کدام کلمات جابه‌جا شده‌اند، معنای یکسان یا کاملاً متفاوتی می‌توانند داشته باشند (برای نمونه، جابه‌جایی نقش فاعل و مفعول در یک جمله، معنا را کاملاً تغییر می‌دهد، درحالی‌که جابه‌جایی یک قید زمانی ممکن است معنا را دست‌نخورده باقی بگذارد)؛ دوم، همین فرایند با ترجمهٔ رفت‌وبرگشتی\LTRfootnote{Back-translation}
(ترجمهٔ جمله به زبانی دیگر و بازگرداندن آن به زبان اصلی) تکرار می‌شود تا تنوع ساختاری بیشتری حاصل شود. برچسب نهایی هر جفت (هم‌معنی/ناهم‌معنی) توسط ارزیاب‌های انسانی تعیین شده است. نتیجه، مجموعه‌ای از جفت‌جمله‌هاست که در آن‌ها، برخلاف
\lr{STS}، همپوشانی واژگانی بالا دیگر نشانهٔ قابل‌اتکایی برای هم‌معنی‌بودن نیست.

نویسندگان
\lr{PAWS}
نشان دادند مدل‌های زمان خود--از جمله مدل‌هایی مبتنی بر
\lr{BERT}
که با رویکردهای بالا آموزش دیده بودند و روی مجموعه‌داده‌های متعارف تشخیص بازنویسی به دقت بالایی می‌رسیدند--در تشخیص این جفت‌های دشوار به‌شدت افت می‌کنند~\cite{zhang2019paws}، و نسخهٔ چندزبانهٔ آن،
\lr{PAWS-X}~\cite{yang2019pawsx}، نشان داد این ضعف به زبان انگلیسی هم محدود نیست. این مجموعه‌داده--که مبنای مستقیم روش پیشنهادی این پایان‌نامه است--بنابراین نه صرفاً یک مجموعه‌دادهٔ تشخیص بازنویسی دیگر، بلکه ابزاری تشخیصی برای آزمودن این فرضیه است که آیا بازنمایی یک مدل، واقعاً بر پایهٔ معنا شکل گرفته یا بر پایهٔ همپوشانی سطحی؛ همین ویژگی است که آن را به گزینه‌ای طبیعی برای سنجش کیفیت بازنمایی لایه‌های میانی یک مدل زبانی--موضوع فصل سوم و چهارم این پایان‌نامه--تبدیل می‌کند.

\subsection{سنجش جداسازی: از دقت طبقه‌بندی تا سطح زیر منحنی}
نکتهٔ روش‌شناختی دیگری که در مرور ادبیات تشخیص بازنویسی اهمیت دارد، \emph{چگونگی سنجش} کیفیت جداسازی است. بسیاری از کارهای اولیه در این حوزه، مسئله را به‌صورت یک وظیفهٔ طبقه‌بندی دودویی کامل صورت‌بندی می‌کنند: یک طبقه‌بندی‌کننده (اغلب یک لایهٔ خطی کوچک روی بازنمایی جفت‌جمله) آموزش داده می‌شود تا برچسب هم‌معنی/ناهم‌معنی را پیش‌بینی کند، و سنجهٔ گزارش‌شده، دقت این طبقه‌بندی‌کننده روی دادهٔ آزمون است. چنین رویکردی، برای هدف این پایان‌نامه دو محدودیت دارد: نخست، به یک مرحلهٔ آموزش جداگانه (آموزش طبقه‌بندی‌کننده) نیاز دارد که خود می‌تواند منبع نویز یا بیش‌برازش باشد، به‌ویژه وقتی هدف، مقایسهٔ ۱۶۰ پیکربندی متفاوت است و آموزش یک طبقه‌بندی‌کنندهٔ جداگانه برای هر پیکربندی، هم هزینهٔ محاسباتی و هم منبع ناهمگونی آماری اضافه می‌کند؛ دوم، دقت طبقه‌بندی، خود به یک آستانهٔ تصمیم دلخواه وابسته است که انتخاب آن می‌تواند نتیجه را تحت‌تأثیر قرار دهد.

به همین دلیل، این پایان‌نامه (همسو با رویکردی که در بسیاری از کارهای سنجش کیفیت بازنمایی جمله، از جمله نسخه‌های اولیهٔ
\lr{STS}
و کارهایی مانند
\lr{Sentence-BERT}~\cite{reimers2019sbert}
که شباهت کسینوسی خام را مستقیماً با سنجه‌های هم‌بستگی می‌سنجند، رایج است) به‌جای دقت طبقه‌بندی، از سطح زیر منحنی مشخصهٔ عملکرد گیرنده
(\lr{AUROC})~\cite{fawcett2006roc}
استفاده می‌کند: با در نظرگرفتن فاصلهٔ کسینوسی هر جفت به‌عنوان یک سنجهٔ پیوستهٔ «امتیاز ناهم‌معنی‌بودن»، بدون آموزش هیچ طبقه‌بندی‌کنندهٔ اضافی و بدون انتخاب هیچ آستانه‌ای، می‌توان کیفیت جداسازی دو توزیع (فاصلهٔ جفت‌های هم‌معنی در برابر فاصلهٔ جفت‌های ناهم‌معنی) را مستقیماً از روی خود بازنمایی سنجید.

یک ویژگی آماری مفید
\lr{AUROC}
که سنجه‌های مبتنی بر دقت فاقد آن‌اند، تفسیر احتمالاتی مستقیم آن است: مقدار
\lr{AUROC}
دقیقاً برابر است با احتمال آنکه، در یک انتخاب تصادفی یک جفت هم‌معنی و یک جفت ناهم‌معنی، فاصلهٔ جفت ناهم‌معنی بزرگ‌تر از فاصلهٔ جفت هم‌معنی باشد. این تفسیر، مقایسهٔ مستقیم نتایج در سرتاسر شبکهٔ آزمایشی این پایان‌نامه--که در آن، لایه‌های مختلف، دامنه‌های مطلقاً متفاوتی برای فاصلهٔ کسینوسی دارند--را بدون نیاز به نرمال‌سازی دستی این دامنه‌ها ممکن می‌سازد؛ ویژگی‌ای که برای طرح آزمایش سه‌بُعدی فصل سوم (شرط × $\lambda$ × لایه) اهمیت عملی مستقیم دارد.

\section{یادگیری تقابلی}
\label{sec:lit-contrastive}

\subsection{مبانی و صورت‌بندی کلی}
یادگیری تقابلی\LTRfootnote{Contrastive Learning}
چارچوبی است برای آموختن یک فضای بازنمایی که در آن، نمونه‌های «مشابه» (جفت‌های مثبت) به هم نزدیک و نمونه‌های «نامشابه» (جفت‌های منفی) از هم دور نگاه داشته می‌شوند، بدون آنکه لزوماً برچسب کلاس صریحی در کار باشد. ریشهٔ این ایده به کارهایی مانند
\lr{FaceNet}~\cite{schroff2015facenet}
در بازشناسی چهره بازمی‌گردد، که با یک تابع هزینهٔ سه‌تایی\LTRfootnote{Triplet Loss}
بر روی سه‌تایی‌های
$(x_a, x_p, x_n)$
(به‌ترتیب: لنگر، مثبت، منفی) به شکل
$\max(0,\ d(x_a,x_p) - d(x_a,x_n) + m)$
عمل می‌کند--یعنی فاصلهٔ لنگر تا نمونهٔ مثبت را دست‌کم به‌اندازهٔ یک حاشیهٔ\LTRfootnote{Margin}
ثابت
$m$
کمتر از فاصلهٔ لنگر تا نمونهٔ منفی نگاه می‌دارد--بازنمایی چهره‌ها را چنان می‌آموزد که فاصلهٔ اقلیدسی مستقیماً معیار شباهت هویت باشد. تابع هزینهٔ ناوردایی این پایان‌نامه (بخش~\ref{sec:loss}) نیز، از نظر ساختاری، بر پایهٔ همین ایده--مقایسهٔ فاصلهٔ جفت «مثبت» با فاصلهٔ جفت «منفی»--بنا شده، هرچند به‌جای یک حاشیهٔ ثابت و دستی، از یک مقیاس و بایاس آموختنی در دل یک تابع هزینهٔ آنتروپی متقاطع دودویی استفاده می‌کند.

صورت‌بندی متداول‌تر امروزی، مبتنی بر تابع هزینهٔ
\lr{InfoNCE}~\cite{oord2018cpc}
است:
\begin{equation}
\mathcal{L}_{\mathrm{InfoNCE}} = -\log \frac{\exp(\mathrm{sim}(x_a, x_p)/\tau)}{\exp(\mathrm{sim}(x_a, x_p)/\tau) + \sum_{k=1}^{K}\exp(\mathrm{sim}(x_a, x_n^{(k)})/\tau)}
\end{equation}
که در آن یک جفت مثبت را در برابر
$K$
نمونهٔ منفی (نه فقط یکی) قرار می‌دهد و آن را به‌صورت یک وظیفهٔ طبقه‌بندی softmax صورت‌بندی می‌کند؛ پارامتر دما\LTRfootnote{Temperature}ی
$\tau$
تندی توزیع حاصل را کنترل می‌کند: مقدار کوچک
$\tau$، تابع هزینه را بیشتر متمرکز بر سخت‌ترین نمونه‌های منفی (آن‌هایی که شباهت بالایی با لنگر دارند) می‌کند، درحالی‌که مقدار بزرگ‌تر آن، وزن یکنواخت‌تری به تمام نمونه‌های منفی می‌دهد؛ این پارامتر، در عمل، نقشی مشابه ضریب مقیاس آموختنی
$\gamma$
در تابع هزینهٔ این پایان‌نامه (بخش~\ref{sec:loss}) ایفا می‌کند، با این تفاوت که در
\lr{InfoNCE}
معمولاً
$\tau$
یک ابرپارامتر ثابت و از پیش‌تعیین‌شده است، درحالی‌که در این پایان‌نامه، مقیاس متناظر، خود در حین آموزش یاد گرفته می‌شود.
\lr{SimCLR}~\cite{chen2020simclr}
در حوزهٔ بینایی ماشین (با استفاده از دو تبدیل تصادفی متفاوت از یک تصویر به‌عنوان جفت مثبت) و
\lr{Supervised Contrastive Learning}~\cite{khosla2020supcon}
(که برچسب کلاس را برای تعیین جفت‌های مثبت/منفی به کار می‌گیرد، به‌جای تکیه بر تبدیل‌های خودکار) نمونه‌های شناخته‌شدهٔ این رویکردند.

\subsection{انتخاب نمونه‌های منفی و نقش نمونه‌های منفی سخت}
\label{sec:lit-hardneg}
کیفیت فضای بازنمایی حاصل از یادگیری تقابلی، به‌شدت به \emph{چگونگی انتخاب نمونه‌های منفی} وابسته است. اگر نمونه‌های منفی به‌سادگی از هر نقطهٔ دیگر فضای داده انتخاب شوند، مدل می‌تواند به‌سرعت یاد بگیرد آن‌ها را با نشانه‌های سطحی و آسان (برای نمونه، طول جمله یا موضوع کلی) از نمونهٔ لنگر تفکیک کند، بدون آنکه مجبور شود به تفاوت‌های ظریف‌تر توجه کند؛ این پدیده در ادبیات با عنوان فروپاشی نمونه‌های منفی آسان\LTRfootnote{Easy Negative Collapse}
شناخته می‌شود. برای مقابله با آن، بسیاری از کارها بر اهمیت نمونه‌های منفی \emph{سخت}\LTRfootnote{Hard Negatives}--نمونه‌هایی که از نظر سطحی به نمونهٔ لنگر شبیه‌اند اما از نظر برچسب هدف، متفاوت‌اند--تأکید کرده‌اند؛ برای نمونه، نسخهٔ نظارت‌شدهٔ
\lr{SimCSE}~\cite{gao2021simcse}
به‌طور صریح از جفت‌های تناقض منطقی (که از نظر موضوع و واژگان به جملهٔ لنگر نزدیک‌اند اما معنایی متضاد دارند) به‌عنوان نمونهٔ منفی سخت استفاده می‌کند.

از این زاویه، مجموعه‌دادهٔ
\lr{PAWS}
را می‌توان نه صرفاً یک مجموعه‌دادهٔ ارزیابی، بلکه یک منبع طبیعی و از پیش‌آماده از نمونه‌های منفی سخت (جفت‌های ناهم‌معنی با همپوشانی واژگانی بسیار بالا) دانست؛ این دقیقاً همان ویژگی‌ای است که این پایان‌نامه از آن برای ساختن جزء «دافعه‌ای»\LTRfootnote{Repulsive}
تابع هزینهٔ ناوردایی (بخش~\ref{sec:loss}) بهره می‌برد. بدون چنین نمونه‌های منفی سختی، تابع هزینهٔ ناوردایی می‌توانست به‌سادگی با فاصله‌گرفتن کلی بازنمایی هر دو جمله از یکدیگر (صرف‌نظر از هم‌معنی‌بودن یا نبودن آن‌ها) کمینه شود؛ وجود جفت‌های ناهم‌معنی لفظاً مشابه، مدل را وامی‌دارد تفاوتی ظریف‌تر و معنایی‌تر بیاموزد.

\subsection{جایگاه این پژوهش در ادبیات یادگیری تقابلی}
تابع هزینهٔ ناوردایی‌ای که در این پایان‌نامه به کار رفته، از نظر مفهومی در همین خانواده جای می‌گیرد--جفت هم‌معنی به‌عنوان جفت مثبت و جفت ناهم‌معنی به‌عنوان جفت منفی سخت--اما دو تفاوت مهم با اکثر ادبیات یادگیری تقابلی دارد. نخست، از نظر صورت‌بندی ریاضی، به‌جای یک وظیفهٔ softmax چندکلاسه (یک مثبت در برابر
$K$
منفی، مطابق
\lr{InfoNCE})، از یک تابع هزینهٔ آنتروپی متقاطع دودویی با مقیاس و بایاس آموختنی استفاده شده--صورت‌بندی‌ای نزدیک‌تر به تابع هزینهٔ سه‌تایی، اما بدون حاشیهٔ دستی‌تنظیم‌شده. دوم، و مهم‌تر، تفاوت این پژوهش با اکثر ادبیات یادگیری تقابلی، \emph{غایت} استفاده از این تابع هزینه است: در
\lr{SimCLR}، \lr{SimCSE}
و مشابه آن‌ها، تابع هزینهٔ تقابلی، تنها یا اصلی‌ترین سیگنال آموزشی برای ساختن یک فضای بازنمایی \emph{از پایه} (یا دست‌کم، برای ریزتنظیم کامل مدل روی این هدف به‌تنهایی) است؛ در این پایان‌نامه، این تابع هزینه یک سیگنال نسبتاً کوچک و مکمل--همراه با یک هزینهٔ زبان‌مدل‌سازی موازی و یک منظم‌ساز بازسازی--است که برای شکل‌دادن \emph{دوبارهٔ} بازنمایی یک مدل از پیش‌آموزش‌دیده، بدون آسیب‌زدن به توانایی‌های عمومی زبانی آن، به کار می‌رود. سومین تفاوت، که موضوع بخش بعدی است، فضایی است که این شباهت کسینوسی در آن محاسبه می‌شود: نه فضای خام بازنمایی مدل، بلکه فضای نهان یک خودرمزگذار پراکنده.

\subsection{ریسک فروپاشی بازنمایی}
یک هشدار شناخته‌شده در ادبیات یادگیری تقابلی، پدیدهٔ فروپاشی بازنمایی\LTRfootnote{Representational Collapse}
است: اگر تابع هزینه صرفاً بر نزدیک‌کردن جفت‌های مثبت متمرکز باشد، بدون هیچ قید متعادل‌کننده‌ای برای دورنگاه‌داشتن نمونه‌های نامرتبط از هم، ساده‌ترین راه‌حلی که بهینه‌سازی می‌تواند بیابد، نگاشت \emph{تمام} ورودی‌ها به یک نقطهٔ ثابت واحد است--راه‌حلی که تابع هزینهٔ مثبت‌محور را دقیقاً صفر می‌کند اما بازنمایی حاصل را کاملاً بی‌فایده می‌سازد، چرا که دیگر هیچ اطلاعاتی دربارهٔ تفاوت ورودی‌ها حمل نمی‌کند. روش‌هایی مانند
\lr{SimCLR}~\cite{chen2020simclr}
و
\lr{InfoNCE}~\cite{oord2018cpc}
این ریسک را با گنجاندن صریح نمونه‌های منفی در همان تابع هزینه (که هرگونه نزدیکی بیش‌ازحد جفت‌های نامرتبط را جریمه می‌کند) کنترل می‌کنند؛ برخی روش‌های جدیدتر بدون نمونهٔ منفی (مانند
\lr{BYOL}
یا
\lr{SimSiam}) از راهکارهای معماری جایگزین (مانند عدم‌تقارن در پس‌انتشار گرادیان میان دو شاخه) برای پیشگیری از فروپاشی استفاده می‌کنند، بدون نیاز به نمونهٔ منفی صریح.

تابع هزینهٔ ناوردایی در این پایان‌نامه (بخش~\ref{sec:loss})، از خانوادهٔ نخست است: جزء دافعه‌ای (جفت‌های ناهم‌معنی
\lr{PAWS}) به همان اندازهٔ جزء جاذبه‌ای، صریحاً در تابع هزینه حضور دارد. این انتخاب، از یک منظر خاص این پایان‌نامه، حتی حیاتی‌تر از نقش متعارفش در جلوگیری از فروپاشی کامل بازنمایی است: چون جفت‌های ناهم‌معنی
\lr{PAWS}
به‌طور خاص، نمونه‌های منفی \emph{سخت} هستند (بخش~\ref{sec:lit-hardneg})، وجودشان در تابع هزینه نه‌تنها از فروپاشی کامل بازنمایی جلوگیری می‌کند، بلکه دقیقاً همان تمایز ظریفی را که هدف اصلی این پایان‌نامه است--تفکیک هم‌معنی از ناهم‌معنی در حضور همپوشانی واژگانی بالا--به‌طور مستقیم در بهینه‌سازی می‌گنجاند؛ بدون این جزء، هیچ تضمینی وجود نداشت که تنها جزء جاذبه‌ای، به تفکیک معنایی--به‌جای فروپاشی ساده--بینجامد.

\section{تفسیرپذیری مکانیستی و خودرمزگذارهای پراکنده}
\label{sec:lit-interpretability}

\subsection{برهم‌نهی و چندمعنایی}
\label{sec:lit-superposition}
حوزهٔ تفسیرپذیری مکانیستی\LTRfootnote{Mechanistic Interpretability}
به دنبال فهم آن است که یک شبکهٔ عصبی آموزش‌دیده، محاسبات درونی خود را با چه سازوکاری انجام می‌دهد؛ نه صرفاً اینکه چه ورودی به چه خروجی می‌انجامد.
\lr{Elhage}
و همکاران~\cite{elhage2021framework}
با ارائهٔ یک چارچوب ریاضی برای تحلیل مدارهای\LTRfootnote{Circuits}
درونی ترنسفورمر--از جمله همین دیدگاه جمعی/انباشتی جریان باقیمانده که در بخش~\ref{sec:lit-representation} معرفی شد--این حوزه را در بستر مدل‌های زبانی مدرن صورت‌بندی کردند.

یکی از یافته‌های مهم این خط پژوهشی، پدیدهٔ برهم‌نهی\LTRfootnote{Superposition}
است.
\lr{Elhage}
و همکاران~\cite{elhage2022superposition}
با استفاده از مدل‌های اسباب‌بازی\LTRfootnote{Toy Models}
کوچک--شبکه‌هایی که مصنوعاً وادار شده‌اند تعداد معینی «ویژگی» ساختگی را در ابعادی کمتر از تعداد آن ویژگی‌ها فشرده کنند--نشان دادند وقتی تعداد مفاهیم بالقوه‌ای که یک شبکه باید نمایش دهد، از تعداد ابعاد موجود در آن لایه بیشتر باشد (وضعیتی که در مدل‌های زبانی واقعی، با میلیون‌ها مفهوم ممکن در برابر چند صد یا چند هزار بُعد، به‌شدت محتمل است)، شبکه--به‌جای تخصیص یک بُعد متعامد به هر مفهوم، که ظرفیت محدود فضا اجازه نمی‌دهد--می‌تواند با قراردادن چند مفهوم به‌صورت همپوشان (نه متعامد، بلکه با زاویه‌ای اندک نسبت به هم) در همان چند بُعد مشترک، بیش از تعداد ابعاد خود، مفهوم نمایش دهد. این راهبرد، تنها زمانی به‌صرفه است که مفاهیم مورد نظر به‌ندرت هم‌زمان فعال شوند (که در دادهٔ زبانی طبیعی، فرض معقولی است: بیشتر مفاهیم ممکن، در بیشتر جمله‌ها اصلاً حضور ندارند)، چرا که برهم‌نهی به بهای مقداری تداخل و اختلال میان مفاهیم هم‌زمان‌فعال تمام می‌شود.

نتیجهٔ مستقیم این پدیده، \emph{چندمعنایی}\LTRfootnote{Polysemanticity}
است: یک نورون منفرد در شبکه، می‌تواند در واکنش به چند مفهوم کاملاً نامرتبط با هم فعال شود (برای نمونه، هم در واکنش به یک الگوی نحوی و هم در واکنش به یک موضوع محتوایی کاملاً بی‌ربط)، و همین امر، تفسیر مستقیم فعال‌سازی تک‌تک نورون‌ها را دشوار می‌سازد؛ نمی‌توان به یک نورون اشاره کرد و گفت «این نورون نمایانگر فلان مفهوم است»، چون همان نورون در بافت‌های دیگر، نمایانگر مفهومی کاملاً متفاوت است. این دقیقاً همان مشکلی است که در بخش~\ref{sec:overview} به‌عنوان دلیل پرهیز از اعمال مستقیم تابع هزینهٔ ناوردایی بر فضای خام جریان باقیمانده مطرح شد: جابه‌جاکردن یک بُعد برهم‌نهی‌شده در یک جهت--برای مثال، برای دورترکردن بازنمایی دو جملهٔ ناهم‌معنی--می‌تواند چندین مفهوم نامرتبط دیگر را که در همان بُعد یا در ابعاد نزدیک به آن جای گرفته‌اند، نیز به‌طور غیرقابل‌پیش‌بینی جابه‌جا کند؛ اثری جانبی که در فضایی غیربرهم‌نهیده، به‌طور ذاتی وجود نخواهد داشت.

\subsection{خودرمزگذارهای پراکنده به‌عنوان راه‌حل}
خودرمزگذارهای پراکنده\LTRfootnote{Sparse Autoencoders, SAEs}
به‌عنوان ابزاری برای گشودن این برهم‌نهی معرفی شدند. ایدهٔ بنیادی، وام‌گرفته از حوزهٔ یادگیری فرهنگ‌لغت\LTRfootnote{Dictionary Learning}
در پردازش سیگنال، این است که اگر بردار فعال‌سازی
$d$
بُعدی یک لایه، در واقع ترکیب خطی پراکنده‌ای از تعداد بسیار بیشتری «جهت ویژگی» زیرین
$m \gg d$
باشد (که پدیدهٔ برهم‌نهی، دقیقاً چنین وضعیتی را پیش‌بینی می‌کند)، آنگاه می‌توان با آموزش یک شبکهٔ رمزگذار--رمزگشای کم‌عمق که فعال‌سازی را به یک کد
$m$
بُعدی پراکنده (اکثر مؤلفه‌ها صفر) می‌برد و سپس تلاش می‌کند فعال‌سازی اصلی را از روی همین کد بازسازی کند، این
$m$
جهت زیرین را بازیابی کرد. تابع هزینهٔ متداول آموزش چنین خودرمزگذاری، به‌طور کلاسیک به شکل
\begin{equation}
\mathcal{L}_{\mathrm{SAE}} = \lVert x - \mathrm{Dec}(\mathrm{Enc}(x)) \rVert_2^2 + \beta \lVert \mathrm{Enc}(x) \rVert_1
\end{equation}
صورت‌بندی می‌شود: جزء نخست، خطای بازسازی استاندارد یک خودرمزگذار، و جزء دوم، جریمه‌ای بر پایهٔ نُرم
$\ell_1$
کد پنهان که--به‌دلیل خاصیت نُرم
$\ell_1$
در ترغیب کدهای دقیقاً صفر (برخلاف نُرم
$\ell_2$
که مقادیر را کوچک اما غیرصفر نگاه می‌دارد)--پراکندگی کد حاصل را تشویق می‌کند؛ ضریب
$\beta$
مصالحه‌ای میان وفاداری بازسازی و درجهٔ پراکندگی را کنترل می‌کند. تابع بخش~\ref{sec:sae-drift} این پایان‌نامه (رابطهٔ~\ref{eq:recon})، از نظر جزء بازسازی، مستقیماً بر همین پایه بنا شده، با این تفاوت که در آن‌جا، به‌جای آموزش خود خودرمزگذار، همین خطای بازسازی به‌عنوان یک سیگنال نظارتی برای \emph{مدل زبانی} (نه خودرمزگذار) به کار می‌رود.

کار
\lr{Cunningham}
و همکاران~\cite{cunningham2023sae}
و هم‌زمان،
\lr{Towards Monosemanticity}
از آنتروپیک~\cite{bricken2023monosemanticity}، نشان دادند این رویکرد--با آموزش یک خودرمزگذار پراکندهٔ کم‌عمق روی فعال‌سازی‌های یک مدل زبانی از پیش‌آموزش‌دیده و \emph{ثابت}--به‌طور سیستماتیک ویژگی‌هایی استخراج می‌کند که به‌مراتب تفسیرپذیرتر از نورون‌های خام‌اند، از ویژگی‌های نحوی و واژگانی ساده (مانند حضور یک کلمهٔ خاص در یک بافت نحوی مشخص) تا مفاهیم انتزاعی‌تر (مانند لحن یا موضوع کلی متن).
\lr{Scaling Monosemanticity}~\cite{templeton2024scaling}
این رویکرد را به مقیاس یک مدل زبانی بزرگ در سطح تولید
(\lr{Claude 3 Sonnet}) گسترش داد و نشان داد ویژگی‌های استخراج‌شده می‌توانند مفاهیمی به‌مراتب انتزاعی‌تر، چندزبانه، و حتی چندوجهی (فعال در واکنش به هم متن و هم تصویر یک مفهوم) را نمایش دهند؛ این کار همچنین نشان داد این ویژگی‌ها، صرفاً همبستگی آماری نیستند، بلکه با دستکاری مقدار آن‌ها (بخش~\ref{sec:lit-steering}) می‌توان رفتار خروجی مدل را در جهت مفهوم متناظر تغییر داد--شاهدی بر نقش علّی این ویژگی‌ها در محاسبات مدل.

از نظر فنی، کار اصلی این پژوهش‌ها--و همچنین کتابخانهٔ
\lr{SAELens}~\cite{bloom2024saelens}
که این پایان‌نامه برای بارگذاری خودرمزگذارهای از پیش‌آموزش‌دیده از آن استفاده کرده--بر پایهٔ همان تابع هزینهٔ بازسازی + جریمهٔ
$\ell_1$
است. توسعه‌های معماری بعدی، تلاش سیستماتیکی برای بهبود مصالحهٔ ذاتی این تابع هزینه--میان پراکندگی کد، وفاداری بازسازی، و جلوگیری از «ویژگی‌های مرده»\LTRfootnote{Dead Features}
(ابعادی از کد که هرگز یا تقریباً هرگز فعال نمی‌شوند و در نتیجه هیچ اطلاعاتی حمل نمی‌کنند)--بوده‌اند. خودرمزگذارهای
\lr{TopK}~\cite{gao2024topk}
به‌جای جریمهٔ نرم
$\ell_1$، مستقیماً و به‌صراحت تنها
$k$
فعال‌ترین بُعد کد را نگه می‌دارند و بقیه را صفر می‌کنند؛ این رویکرد، پراکندگی کد را به‌طور دقیق و از پیش‌تعیین‌شده تضمین می‌کند، بدون آنکه نیاز به تنظیم ظریف ضریب جریمهٔ
$\ell_1$
باشد. خودرمزگذارهای
\lr{JumpReLU}~\cite{rajamanoharan2024jumprelu}
رویکرد دیگری در پیش گرفتند: با جایگزین‌کردن تابع فعال‌سازی
\lr{ReLU}
با یک تابع آستانه‌گذاری آموختنی (که برای هر بُعد کد، آستانهٔ متفاوتی می‌آموزد)، مستقیماً مصالحهٔ میان پراکندگی و خطای بازسازی را بهینه می‌کنند، بدون کوچک‌ترشدن سیستماتیک دامنهٔ ویژگی‌های واقعاً فعال--اثر جانبی ناخواسته‌ای که جریمهٔ
$\ell_1$
علاوه بر تشویق پراکندگی، به‌طور طبیعی به همراه دارد.

این پایان‌نامه، به‌جای آموزش یک خودرمزگذار جدید یا مقایسهٔ معماری‌های مختلف آن، از یک خودرمزگذار از پیش‌آموزش‌دیده و منتشرشده (با معماری استاندارد بازسازی+
$\ell_1$) به‌عنوان جزئی ثابت و آماده استفاده می‌کند؛ سهم این پایان‌نامه، نه در بهبود خود خودرمزگذار، بلکه در یافتن کاربرد تازه‌ای برای آن است (بخش~\ref{sec:lit-steering}).

\subsection{فراتر از تحلیل: هدایت و مهندسی بازنمایی}
\label{sec:lit-steering}
کاربرد غالب خودرمزگذارهای پراکنده، همچنان تحلیل \emph{پس‌ازآموزش} یک مدل ثابت است: مدل آموزش داده می‌شود، سپس خودرمزگذار روی فعال‌سازی‌های آن آموزش می‌بیند، و سرانجام ویژگی‌های حاصل بررسی و نام‌گذاری می‌شوند. اما این حوزه به‌تدریج به کاربردهای \emph{مداخله‌گر} نیز گسترش یافته است.
\lr{Turner}
و همکاران~\cite{turner2023actadd}
نشان دادند می‌توان با افزودن یک بردار جهت ثابت (برآمده از تفاوت فعال‌سازی دو دستهٔ متن متضاد، مانند متن «شاد» در برابر «غمگین») به فعال‌سازی مدل در زمان استنتاج، رفتار خروجی مدل را--بدون هیچ‌گونه ریزتنظیمی و تنها با یک عملیات برداری ساده در زمان اجرا--در جهت مورد نظر «هدایت»\LTRfootnote{Steering}
کرد.
\lr{Scaling Monosemanticity}~\cite{templeton2024scaling}
نیز نشان داد وقتی این بردارهای جهت از میان ویژگی‌های یک خودرمزگذار پراکنده (به‌جای جهت‌های دستی‌ساخته از تفاضل دو دسته متن) استخراج شوند، هدایت حاصل می‌تواند دقیق‌تر و موضعی‌تر باشد، چون هر ویژگی، به‌لطف پراکندگی کد، تداخل کمتری با مفاهیم نامرتبط دارد.
\lr{Zou}
و همکاران، در چارچوبی که «مهندسی بازنمایی»\LTRfootnote{Representation Engineering}~\cite{zou2023repe}
می‌نامند، این ایده را به یک رویکرد کلی‌تر بالا-به-پایین برای شفافیت هوش مصنوعی--شناسایی جهت‌های مفهومی (نه لزوماً از طریق خودرمزگذار، بلکه با روش‌های ساده‌تر آماری) و مداخله در آن‌ها، به‌جای تحلیل مدارهای پایین-به-بالا و مدارمحور--تعمیم داده‌اند.

همهٔ کاربردهای بالا--چه هدایت مبتنی بر تفاضل فعال‌سازی و چه هدایت مبتنی بر ویژگی خودرمزگذار--یک ویژگی مشترک دارند: خودرمزگذار (یا بردار جهت)، پس از آموزش کامل مدل و به‌عنوان یک لایهٔ مداخلهٔ \emph{زمان استنتاج} به کار می‌رود؛ وزن‌های خود مدل زبانی دست‌نخورده باقی می‌مانند، و هدایت باید در هر بار فراخوانی مدل، دوباره اعمال شود. رویکرد این پایان‌نامه، از این جهت، ماهیتاً متفاوت است: خودرمزگذار نه در زمان استنتاج، بلکه در \emph{زمان آموزش} و نه به‌عنوان مداخله‌گر بلکه به‌عنوان \emph{فضای محاسبهٔ تابع هزینه} به کار می‌رود؛ نتیجهٔ نهایی، تغییری ماندگار در خود وزن‌های مدل (از طریق لورا) است، نه یک بردار افزودنی جداگانه که در زمان استنتاج باید اعمال شود--با این پیامد که، برخلاف هدایت زمان استنتاج، خودرمزگذار پس از پایان آموزش دیگر برای بهره‌گیری از مدل ریزتنظیم‌شده لازم نیست. تا آنجا که در چارچوب این پایان‌نامه بررسی شده، این کاربرد--استفاده از یک خودرمزگذار ثابت به‌عنوان عدسی یک تابع هزینهٔ آموزشی--در ادبیات موجود کمتر مورد بررسی مستقیم قرار گرفته است.

\subsection{نقد و محدودیت‌های کاربردی}
شایان ذکر است که سودمندی عملی خودرمزگذارهای پراکنده--حتی در همان کاربردهای مرسوم‌تر--بدون چالش نیست، و این پایان‌نامه از پذیرش بی‌قیدوشرط برتری آن‌ها پرهیز می‌کند.
\lr{Kantamneni}
و همکاران~\cite{kantamneni2025useful}
در بررسی کاربرد خودرمزگذارها برای کاوش\LTRfootnote{Probing}
مفاهیم--یعنی آموختن یک طبقه‌بندی‌کنندهٔ ساده روی کد خودرمزگذار برای پیش‌بینی یک ویژگی مورد نظر--نشان دادند در بسیاری از تنظیمات، یک طبقه‌بندی‌کنندهٔ خطی مستقیم روی فعال‌سازی خام (بدون هیچ خودرمزگذاری) می‌تواند عملکردی هم‌تراز یا حتی بهتر داشته باشد؛ به بیان دیگر، پراکنده‌ترکردن بازنمایی، همیشه یک مسئلهٔ پایین‌دستی نظارت‌شده را ساده‌تر نمی‌کند. به‌طور مشابه،
\lr{Wu}
و همکاران در چارچوب
\lr{AxBench}~\cite{wu2025axbench}
در بررسی کاربرد خودرمزگذارها برای هدایت رفتار مدل (بخش~\ref{sec:lit-steering})، نشان دادند در مقایسه با پایه‌های ساده--مانند تفاضل‌میانگین فعال‌سازی میان دو دستهٔ متن متضاد، بدون هیچ خودرمزگذاری--بردارهای برآمده از خودرمزگذار، همیشه برتری قابل‌توجهی در کیفیت هدایت نشان نمی‌دهند و گاه حتی عملکرد ضعیف‌تری دارند.

این یافته‌ها هشداری روش‌شناختی برای این پایان‌نامه‌اند: نمی‌توان صرفاً فرض کرد فضای خودرمزگذار «به‌طور ذاتی» و بدون هیچ آزمونی، برای هر کاربردی--از جمله کاربرد تازهٔ این پایان‌نامه، یعنی محاسبهٔ یک تابع هزینهٔ آموزشی--بهتر از فضای خام است؛ برتری آن باید در هر کاربرد مشخص، به‌طور تجربی و در برابر یک مبنای مقایسه‌ای مبتنی بر فضای خام سنجیده شود. این پایان‌نامه دقیقاً همین رویکرد را، با گنجاندن شرط کنترلی
\lr{C2} (اعمال همان تابع هزینه، اما در فضای خام؛ بخش~\ref{sec:conditions})، در پیش گرفته است؛ نتایج فصل چهارم (بخش~\ref{sec:res-condition}) نشان می‌دهند که، دست‌کم در این کاربرد خاص و بر روی این مدل، فضای خودرمزگذار واقعاً برتری قابل‌سنجش دارد--اما این نتیجه، خروجی یک آزمایش کنترل‌شده است، نه یک فرض پیشینی.

\section{ریزتنظیم کارآمد پارامتر و حفظ رفتار مدل}
\label{sec:lit-peft}

\subsection{روش‌های ریزتنظیم کارآمد پارامتر}
با بزرگ‌ترشدن مدل‌های زبانی، ریزتنظیم کامل تمام پارامترها برای هر وظیفهٔ جدید، هم از نظر محاسباتی و هم از نظر حافظه (چرا که بهینه‌سازهایی مانند
\lr{Adam}
به ازای هر پارامتر، چند مقدار کمکی اضافی نگه می‌دارند)، پرهزینه شده است. روش‌های ریزتنظیم کارآمد پارامتر\LTRfootnote{Parameter-Efficient Fine-Tuning, PEFT}
با این هدف طراحی شده‌اند که تنها بخش کوچکی از پارامترها (یا پارامترهای افزودنی جدید) آموزش داده شوند، درحالی‌که بیشتر وزن‌های از پیش‌آموزش‌دیده ثابت باقی می‌مانند.
\lr{Adapters}~\cite{houlsby2019adapters}
از نخستین رویکردهای این خانواده بودند: افزودن لایه‌های کوچک میانی (معمولاً یک گلوگاه دوبُعدی کاهشی-افزایشی) به هر بلوک ترنسفورمر، به‌گونه‌ای که تنها همین لایه‌های افزوده‌شده آموزش داده شوند.
\lr{Prefix-Tuning}~\cite{li2021prefixtuning}
و
\lr{Prompt Tuning}~\cite{lester2021prompttuning}
رویکرد دیگری در پیش گرفتند: به‌جای تغییر معماری درونی مدل، بردارهای پیوستهٔ آموختنی («پیشوند» یا «اعلان») به ورودی مدل (یا به حالت‌های کلید-مقدار توجه) افزوده می‌شوند، و مدل اصلی کاملاً دست‌نخورده باقی می‌ماند.

لورا~\cite{hu2021lora}--روشی که این پایان‌نامه از آن استفاده کرده--رویکرد سومی را پیش گرفت: به‌جای افزودن لایه یا توکن جدید، یک به‌روزرسانی کم‌رتبه مستقیماً به ماتریس‌های وزن موجود افزوده می‌شود (جزئیات ریاضی در بخش~\ref{sec:base-model}). مزیت عملی این رویکرد نسبت به دو رویکرد پیشین دوگانه است: نخست، پس از آموزش، می‌توان به‌روزرسانی کم‌رتبه را مستقیماً در وزن اصلی ادغام کرد و هیچ سربار محاسباتی‌ای--نه در طول دنباله (برخلاف
\lr{Prefix-Tuning}
که طول مؤثر دنباله را افزایش می‌دهد) و نه در ساختار شبکه (برخلاف
\lr{Adapters})--در زمان استنتاج باقی نمی‌ماند؛ دوم، فضای پارامترهای قابل‌آموزش (ماتریس‌های
$A$
و
$B$) به‌طور طبیعی با رتبهٔ
$r$
مقیاس‌پذیر است، که امکان کنترل دقیق مصالحهٔ میان ظرفیت تطبیق و صرفه‌جویی محاسباتی را فراهم می‌کند.

از منظر این پایان‌نامه، معیار انتخاب میان این رویکردهای مختلف
\lr{PEFT}
تنها کارآمدی پارامتر نبود، بلکه سازگاری با طرح آزمایش شبکه‌ای (بخش~\ref{sec:ablation}) نیز اهمیت داشت: از آنجا که هر یک از ۱۶۰ پیکربندی باید به‌طور کاملاً مستقل و از ابتدا آموزش می‌دید، روشی که پس از آموزش، بار اضافه‌ای (نه در معماری و نه در طول دنباله) به مدل نهایی تحمیل نکند، برای مدیریت‌پذیری عملی این حجم از اجراها ترجیح داده شد؛ لورا، با قابلیت ادغام‌شدن کامل خود در وزن اصلی پس از آموزش، این ویژگی را به‌طور طبیعی برآورده می‌کند. افزون‌بر این، از آنجا که تابع هزینهٔ ناوردایی (بخش~\ref{sec:loss}) بر روی فعال‌سازی جریان باقیمانده در یک لایهٔ \emph{مشخص} تعریف می‌شود، روشی که مستقیماً وزن‌های همان لایه‌ها را تطبیق دهد (مانند لورا، که می‌تواند به‌صورت گزینشی تنها به ماتریس‌های وزن لایه‌های مورد نظر اعمال شود)، از نظر مفهومی سرراست‌تر از روشی است که طول دنباله یا معماری شبکه را تغییر می‌دهد.

\subsection{حفظ رفتار مدل در حین ریزتنظیم}
ریزتنظیم کارآمد پارامتر، مسئلهٔ حجم پارامترهای آموزش‌دیده را حل می‌کند، اما مسئلهٔ جداگانه‌ای را حل‌نشده باقی می‌گذارد: حتی با تعداد کمی پارامتر قابل‌آموزش، ریزتنظیم متمرکز بر یک وظیفهٔ باریک می‌تواند توانایی‌های عمومی‌تر مدل را تضعیف کند؛ پدیده‌ای که در ادبیات یادگیری پیوسته با عنوان فراموشی فاجعه‌بار\LTRfootnote{Catastrophic Forgetting}
شناخته می‌شود.
\lr{Kirkpatrick}
و همکاران~\cite{kirkpatrick2017ewc}
یکی از نخستین راه‌حل‌های سیستماتیک برای این مسئله را با روش تحکیم وزن کشسان\LTRfootnote{Elastic Weight Consolidation, EWC}
ارائه کردند: پارامترهایی که برای وظیفهٔ قبلی اهمیت بیشتری داشته‌اند (بر اساس اطلاعات فیشر\LTRfootnote{Fisher Information}، که به‌طور تقریبی، حساسیت تابع هزینهٔ وظیفهٔ قبلی را نسبت به هر پارامتر می‌سنجد) با جریمهٔ سنگین‌تری--متناسب با مربع فاصله از مقدار قبلی‌شان--از تغییر زیاد بازداشته می‌شوند؛ پارامترهایی که اهمیت کمتری داشته‌اند، آزادانه‌تر تغییر می‌کنند.

الگوی مشابهی، اما در سطح \emph{رفتار خروجی} مدل به‌جای پارامترهای منفرد، در ادبیات یادگیری تقویتی مبتنی بر بازخورد انسانی\LTRfootnote{Reinforcement Learning from Human Feedback, RLHF}
دیده می‌شود.
\lr{InstructGPT}~\cite{ouyang2022instructgpt}
نشان داد وقتی یک مدل زبانی صرفاً با هدف بیشینه‌سازی یک سیگنال پاداش (اینجا: ترجیح انسانی، خود برآمده از یک مدل پاداش جداگانه) بهینه شود، می‌تواند به شکلی نامطلوب--با یافتن راه‌های میان‌بُری که مدل پاداش را فریب می‌دهند بدون آنکه واقعاً کیفیت خروجی را بهبود دهند--از رفتار زبانی اصیل خود و حتی از حقیقت دور شود.
% پدیده‌ای که گاه با اصطلاح عمومی‌تر «قانون گودهارت»\LTRfootnote{Goodhart's Law}
% (هر سنجه‌ای که به هدف بهینه‌سازی تبدیل شود، دیگر سنجهٔ خوبی نیست) توصیف می‌شود
برای مهار این پدیده، یک جریمهٔ واگرایی کولبک--لایبلر نسبت به سیاست اولیه (پیش از یادگیری تقویتی) به تابع هزینه افزوده شد، دقیقاً برای آنکه مدل در ناحیه‌ای از فضای رفتار بماند که «شبیه به» مدل اصیل باقی بماند، درحالی‌که هم‌زمان در جهت هدف جدید (ترجیح انسانی) نیز پیش می‌رود.

\subsection{اعتبار ابزارهای سنجش ثابت در حین آموزش}
نکته‌ای که در هر دو مثال بالا مشترک است--و مستقیماً به ریسک رانش خودرمزگذار این پایان‌نامه (بخش~\ref{sec:sae-drift}) مرتبط می‌شود--وجود یک «ابزار سنجش ثابت» است که در حین بهینه‌سازی، خودش تغییر نمی‌کند اما مبنای سنجش موفقیت یا محدودیت قرار می‌گیرد: در
\lr{EWC}، این ابزار، مقدار اطلاعات فیشر محاسبه‌شده روی مدل اولیه است؛ در
\lr{RLHF}، این ابزار، هم سیاست اولیه (برای جریمهٔ کولبک--لایبلر) و هم خود مدل پاداش است، که اگر سیاست درحال‌آموزش بیش‌ازحد از توزیع داده‌ای که مدل پاداش روی آن آموزش دیده دور شود، دیگر نمی‌توان به خروجی آن اعتماد کرد. در این پایان‌نامه، این ابزار سنجش ثابت، خودرمزگذار پراکنده است: سنجهٔ اصلی موفقیت (جداسازی هم‌معنی/ناهم‌معنی) تنها زمانی معنادار است که خودرمزگذار همچنان بازسازی‌کنندهٔ باکیفیتی از فعال‌سازی‌های مدل باشد.

منظم‌ساز بازسازی این پایان‌نامه (بخش~\ref{sec:sae-drift})، از نظر نقش مفهومی--نه از نظر فرمول ریاضی--با هر سه رویکرد بالا خویشاوند است: در همهٔ موارد، یک جزء تابع هزینه، مدل را از دورشدن بیش‌ازحد از یک نقطهٔ مرجع ثابت بازمی‌دارد، درحالی‌که جزء دیگر، مدل را در جهت هدف تازه پیش می‌راند. تفاوت کلیدی آن است که در
\lr{EWC}، نقطهٔ مرجع، \emph{پارامترهای} مدل هستند؛ در
\lr{RLHF}، نقطهٔ مرجع، \emph{توزیع خروجی} مدل و اعتبار مدل پاداش است؛ و در این پایان‌نامه، نقطهٔ مرجع، نه پارامتر و نه خروجی، بلکه \emph{قابلیت بازسازی فعال‌سازی‌ها توسط یک ابزار سنجش شخص ثالث} است.% این دوگانگی نقش (خودرمزگذار هم فضای هدف است و هم ابزار سنجش اعتبار خودش) گزینه‌ای است که، تا آنجا که در این پایان‌نامه بررسی شده، در ادبیات منظم‌سازی فراموشی، پیشینهٔ مستقیمی ندارد، و پیامدهای تجربی آن--به‌ویژه محدودیت‌های منظم‌ساز در بازهٔ وزن‌های بالای این هدف--موضوع اصلی بخش~\ref{sec:res-dynamics} از فصل چهارم است.

\section{جمع‌بندی و جایگاه این پژوهش}
\label{sec:lit-summary}
مرور بالا نشان می‌دهد سه ادعای بنیادی این پایان‌نامه، هرکدام به‌طور مستقل در ادبیات موجود پشتوانه دارند: اینکه بازنمایی خام مدل‌های زبانی، حتی وقتی بافتاری است، لزوماً معناداری‌ای در سطح مورد نظر انسان ندارد، و مجموعه‌دادهٔ
\lr{PAWS}
ابزار مناسبی برای آشکارسازی این نبود است (بخش~\ref{sec:lit-representation})؛ اینکه یک تابع هزینهٔ تقابلی، به‌ویژه با تکیه بر نمونه‌های منفی سخت، می‌تواند فضای بازنمایی را در جهت یک سنجهٔ معنایی مشخص شکل دهد (بخش~\ref{sec:lit-contrastive})؛ و اینکه فضای خام فعال‌سازی، به‌دلیل برهم‌نهی ویژگی‌ها، مسیر ایده‌آلی برای این نوع مداخله نیست، درحالی‌که خودرمزگذارهای پراکنده برای گشودن همین برهم‌نهی طراحی شده‌اند--هرچند سودمندی عملی آن‌ها، آن‌گونه که نقدهای اخیر نشان می‌دهند، باید در هر کاربرد به‌طور مستقل آزموده شود، نه مفروض گرفته شود (بخش~\ref{sec:lit-interpretability}).

جدول~\ref{tab:lit-position} این پایان‌نامه را در برابر نزدیک‌ترین خطوط پژوهشی مرورشده، از نظر سه محور کلیدی، جای‌گذاری می‌کند.

\begin{table}[htb]
\centering
\caption{جایگاه این پژوهش نسبت به نزدیک‌ترین خطوط پژوهشی مرورشده.}
\label{tab:lit-position}
\begin{tabular}{p{3.5cm}p{3.3cm}p{3.3cm}p{3.2cm}}
\toprule
 & نقش خودرمزگذار & زمان اعمال & اثر بر وزن‌های مدل \\
\midrule
تفسیرپذیری پس‌ازآموزش (مثل~\cite{bricken2023monosemanticity}) & ابزار تحلیل و کشف & پس از آموزش کامل & بدون اثر \\
هدایت زمان استنتاج (مثل~\cite{templeton2024scaling,turner2023actadd}) & اهرم مداخله & زمان استنتاج، هر بار & بدون اثر \newline (مداخله‌ای جدا) \\
یادگیری تقابلی متعارف (مثل~\cite{gao2021simcse}) & بدون خودرمزگذار \newline (فضای خام یا مدل جدا) & زمان آموزش & مستقیم \\
\textbf{این پایان‌نامه} & فضای هدف تابع هزینه & زمان آموزش & مستقیم \\
\bottomrule
\end{tabular}
\end{table}

آنچه در ادبیات بررسی‌شده یافت نشد--و شکافی است که این پایان‌نامه می‌کوشد آن را پر کند--\emph{ترکیب} این سه ادعا در یک روش واحد است، آن‌گونه که در ستون آخر جدول بالا خلاصه شده: استفاده از خودرمزگذار پراکنده، نه به‌عنوان ابزار تحلیل پس‌ازآموزش و نه به‌عنوان اهرم هدایت زمان استنتاج، بلکه به‌عنوان فضای هدف یک تابع هزینهٔ تقابلی که مستقیماً وزن‌های مدل را--با روشی کارآمد پارامتر و با یک منظم‌ساز صریح برای حفظ اعتبار همان خودرمزگذار (بخش~\ref{sec:lit-peft})--در حین آموزش تغییر می‌دهد. فصل سوم، پیاده‌سازی دقیق این ترکیب را شرح می‌دهد.
